% =====================================================================
% Định dạng báo cáo, theo mẫu báo cáo môn học của trường (file Word cũ):
% A4; lề trên/dưới 3 cm, trái 3.5 cm, phải 2 cm; Times 13 pt; giãn dòng 1.5;
% tên chương 15 pt đậm, căn giữa; chú thích hình dưới hình, bảng trên bảng.
% Biên dịch bằng XeLaTeX (Overleaf: Menu -> Compiler -> XeLaTeX).
%
% Overleaf miễn phí chỉ cho 10 giây mỗi lần biên dịch, nên preamble chỉ nạp gói thật cần:
% toán dùng newtxmath (nhanh hơn nhiều so với unicode-math), không nạp TikZ (sơ đồ vẽ sẵn thành PDF),
% không dùng độ rộng số tự động trong mục lục (nó bắt biên dịch thêm một lượt).
% =====================================================================

% ---------- Font: Times cho chữ (TeX Gyre Termes, đủ dấu tiếng Việt) và cho toán (newtxmath) ----------
\usepackage{amsmath}
\usepackage{newtxmath}
\usepackage[no-math]{fontspec}
\setmainfont{texgyretermes}[
  Extension      = .otf,
  UprightFont    = *-regular,
  BoldFont       = *-bold,
  ItalicFont     = *-italic,
  BoldItalicFont = *-bolditalic]

% ---------- Ngôn ngữ ----------
\usepackage{polyglossia}
\setdefaultlanguage{vietnamese}
\setotherlanguage{english}
\usepackage{etoolbox}
\newcommand{\setreportnames}{%
  \renewcommand{\contentsname}{MỤC LỤC}%
  \renewcommand{\listfigurename}{DANH SÁCH HÌNH ẢNH}%
  \renewcommand{\listtablename}{DANH SÁCH BẢNG}%
  \renewcommand{\bibname}{TÀI LIỆU THAM KHẢO}%
  \renewcommand{\figurename}{Hình}%
  \renewcommand{\tablename}{Bảng}%
  \renewcommand{\chaptername}{CHƯƠNG}%
}
% Đặt cả hai chỗ: cách polyglossia nạp tên mục khác nhau giữa các phiên bản TeX Live
\gappto\captionsvietnamese{\setreportnames}
\AtBeginDocument{\setreportnames}

% ---------- Trang ----------
\usepackage[a4paper, top=3cm, bottom=3cm, left=3.5cm, right=2cm,
            headheight=24pt, headsep=0.6cm, footskip=1.2cm]{geometry}
\usepackage{setspace}
% Word "1.5 lines" với Times ~ 1.5 x 1.15 cỡ chữ; LaTeX mặc định 1.2 cỡ chữ -> hệ số 1.5 x 1.15 / 1.2 ~ 1.43
\setstretch{1.43}
\setlength{\parindent}{1cm}
\setlength{\parskip}{3pt plus 1pt}
\usepackage{indentfirst}
% Tên bộ dữ liệu, tham số (\texttt) không ngắt được: cho phép giãn khoảng trắng thêm thay vì tràn lề
\setlength{\emergencystretch}{3em}

% ---------- Header / footer: tên đề tài có gạch dưới, số trang giữa chân trang ----------
\usepackage[headsepline]{scrlayer-scrpage}
\setlength{\footheight}{24pt}
\clearpairofpagestyles
% Cỡ 9 pt để tên đề tài nằm trên một dòng như header của mẫu
\ihead{\fontsize{9}{11}\selectfont ỨNG DỤNG HỌC SÂU TRONG PHÁT HIỆN, PHÂN LOẠI VÀ THEO DÕI PHƯƠNG TIỆN GIAO THÔNG}
\cfoot*{\pagemark}
\setkomafont{pageheadfoot}{\normalfont}
\setkomafont{pagenumber}{\normalfont}
\renewcommand*{\chapterpagestyle}{scrheadings}

% ---------- Tiêu đề chương / mục ----------
\setkomafont{disposition}{\normalfont\bfseries}
\setkomafont{chapter}{\fontsize{15}{19}\selectfont\bfseries}
\setkomafont{section}{\fontsize{13}{16}\selectfont\bfseries}
\setkomafont{subsection}{\fontsize{13}{16}\selectfont\bfseries}
\setkomafont{subsubsection}{\fontsize{13}{16}\selectfont\bfseries\itshape}
\renewcommand*{\raggedchapter}{\centering}
\renewcommand*{\chapterformat}{\chaptername~\thechapter.\space}
\RedeclareSectionCommand[beforeskip=0pt, afterskip=1.2\baselineskip]{chapter}
\RedeclareSectionCommands[beforeskip=0.8\baselineskip, afterskip=0.4\baselineskip]{section,subsection}
\RedeclareSectionCommand[beforeskip=0.6\baselineskip, afterskip=0.3\baselineskip]{subsubsection}
\setcounter{secnumdepth}{\subsubsectionnumdepth}
\setcounter{tocdepth}{\subsectiontocdepth}

% Mục lục ghi "CHƯƠNG 1. TÊN", danh sách hình/bảng ghi "Hình 1.1 ..."/"Bảng 1.1 ...".
% Độ rộng ô số đặt cố định (vừa "CHƯƠNG 4." và "Hình 4.10") thay cho tính tự động.
\setuptoc{toc}{totoc}  % đưa chính MỤC LỤC vào mục lục như mẫu
\newcommand*{\chapterprefixintoc}[1]{\chaptername~#1\unskip.\space}
\RedeclareSectionCommand[tocentrynumberformat=\chapterprefixintoc, tocnumwidth=2.6cm]{chapter}
\newcommand*{\figureprefixintoc}[1]{\figurename~#1}
\newcommand*{\tableprefixintoc}[1]{\tablename~#1}
\DeclareTOCStyleEntry[entrynumberformat=\figureprefixintoc, numwidth=2.1cm]{tocline}{figure}
\DeclareTOCStyleEntry[entrynumberformat=\tableprefixintoc, numwidth=2.1cm]{tocline}{table}

% ---------- Hình, bảng, chú thích ----------
\usepackage{graphicx}
\graphicspath{{figures/}}
\usepackage{float}
\usepackage[table]{xcolor}
\definecolor{tableheader}{HTML}{5B9BD5}
\usepackage{array, tabularx, longtable, multirow}
\newcolumntype{C}{>{\centering\arraybackslash}X}
\newcolumntype{L}{>{\raggedright\arraybackslash}X}
\renewcommand{\arraystretch}{1.15}
\usepackage{caption}
\usepackage{subcaption}
\DeclareCaptionFont{twelve}{\fontsize{12}{14.5}\selectfont}
\captionsetup{font={twelve,it}, labelsep=space, justification=centering, skip=6pt}
\captionsetup[table]{position=top}
% Hàng tiêu đề bảng: nền xanh, chữ đậm như mẫu
\newcommand{\thead}[1]{\cellcolor{tableheader}\textbf{#1}}

% ---------- Danh sách ----------
\usepackage{enumitem}
\setlist[itemize,1]{label=--, leftmargin=1.5cm, itemsep=1pt, topsep=2pt}
\setlist[itemize,2]{label=+, leftmargin=*, itemsep=0pt}
\setlist[enumerate]{leftmargin=1.5cm, itemsep=1pt, topsep=2pt}

% ---------- Toán ----------
\numberwithin{equation}{chapter}

% ---------- Khung trang bìa (eso-pic nhẹ hơn TikZ và không cần biên dịch thêm lượt) ----------
\usepackage{eso-pic}

% ---------- Khác ----------
% Không dùng hyperref: báo cáo nộp bản in, bỏ liên kết bấm được để mỗi lượt biên dịch nhanh hơn ~0.3 s
\usepackage{url}
\urlstyle{same}
